\documentclass[10pt,a4paper]{article}
\usepackage[T1]{fontenc}
\usepackage[utf8]{inputenc}
\usepackage[top=1.5cm, bottom=1.0cm, left=1.8cm, right=1.8cm]{geometry}
\usepackage{booktabs}
\usepackage{amsmath}
\usepackage{parskip}
\usepackage{microtype}
\usepackage{enumitem}
\usepackage{titlesec}
\usepackage{xcolor}
\setlength{\parskip}{3pt}
\setlength{\parindent}{0pt}
\titleformat{\subsection}{\bfseries\normalsize}{}{0pt}{}[{\vspace{1pt}\hrule\vspace{2pt}}]
\titlespacing{\subsection}{0pt}{6pt}{0pt}
\pagestyle{empty}
\begin{document}

\begin{center}
    {\LARGE \textbf{Decarbonise European Equities Without Giving Up the Benchmark}}\\[0.10cm]
    {\large A Practical, Rules-Based Decarbonisation Framework}
\end{center}
\rule{\linewidth}{1pt}
\vspace{-0.25cm}

Institutional investors face a clear challenge: reduce carbon exposure
without materially weakening the financial profile of a broad European
equity allocation. Using 12 years of out-of-sample data on 633 European
equities, we find that a rules-based decarbonisation strategy can deliver
large carbon reductions at a limited observed return cost.

\subsection{Two Practical Portfolio Solutions}
\begin{enumerate}[leftmargin=*, itemsep=0pt, topsep=2pt]
    \item \textbf{50\% Carbon Cap} ($P^{vw}(0.5)$): stays close to the
    European value-weighted benchmark while halving the benchmark's carbon
    footprint each year.
    \item \textbf{Net-Zero Glide Path} ($P^{vw}(NZ)$): reduces the carbon
    budget by 10\% per year from the 2013 baseline, aligned with the EU
    Climate Transition Benchmark.
\end{enumerate}

\subsection{Headline Results --- Out-of-Sample, 2014--2025}
\begin{center}
\renewcommand{\arraystretch}{1.15}
\colorbox{black!5}{%
\begin{tabular}{lccc}
\toprule
\textbf{Strategy} & \textbf{Geometric Return} & \textbf{Gap vs.\ Benchmark} & \textbf{Carbon Reduction} \\
\midrule
Benchmark $P^{vw}$       & 7.16\%             & ---                 & ---                      \\
50\% Cap $P^{vw}(0.5)$   & 6.78\%             & $-$38 bp            & $-$50\%                  \\
Net-Zero $P^{vw}(NZ)$    & \textbf{6.84\%}    & \textbf{$-$32 bp}   & \textbf{$-$72\% by 2024} \\
\bottomrule
\end{tabular}}
\end{center}
\vspace{-2pt}
\begin{center}
{\footnotesize\textit{Zero risk-free rate. Substantial decarbonisation at a
modest observed return cost; the Net-Zero glide path is slightly cheaper
than the static cap because the benchmark self-decarbonised over the
period.}}
\end{center}

\subsection{Why the Strategy Works}
\begin{itemize}[leftmargin=*, itemsep=0pt, topsep=2pt]
    \item \textbf{Carbon is concentrated.} 10 firms (Holcim, RWE, Rio Tinto,
    Enel, \dots) account for 44\% of the benchmark's WACI in 2013 ---
    selective underweighting captures most of the reduction.
    \item \textbf{The benchmark self-decarbonised.} The European
    value-weighted CF fell by 58\% between 2013 and 2024, which made the
    Net-Zero glide path particularly efficient.
\end{itemize}

\subsection{Risks and Trade-offs}
\begin{itemize}[leftmargin=*, itemsep=0pt, topsep=2pt]
    \item \textbf{Sample dependence:} results reflect European equities over
    2014--2025; carbon-premium effects may differ in other markets or windows.
    \item \textbf{Scope 3 excluded:} the analysis uses Scope 1 + Scope 2
    emissions only.
    \item \textbf{Implementation frictions:} transaction costs, turnover and
    trading constraints to be assessed before deployment.
\end{itemize}

\subsection{Recommendation}
For immediate, measurable decarbonisation, the \textbf{50\% Carbon Cap}
delivers a clear, transparent solution with limited tracking deviation. For
long-term regulatory alignment and the deepest carbon reduction, the
\textbf{Net-Zero Glide Path} is the stronger choice: it ends the horizon
with $-$72\% carbon at a smaller observed return cost than the static cap.
Both approaches are transparent, rules-based, and suitable for investors
seeking lower financed emissions while preserving a benchmark-like European
equity profile.

\vfill
\rule{\linewidth}{0.5pt}
\begin{center}
{\footnotesize \textbf{Group AU} --- HEC Lausanne --- Sustainability-Aware
Asset Management --- May 2026 \quad \textbar \quad Arthur Pillet \quad
Nathan Heimendinger Dreyfus \quad Marc Muller \quad Evan Dejean}
\end{center}

\end{document}
